\begin{python0}
from solutions import *; clear()
\end{python0}
\sectionthree{static\_cast}

The \texttt{\textbf{static\_cast}} is most similar to the basic
typecasting of types which are very similar or only result in loss of
precision, i.e., casting such as
\begin{console}
int i = 42;
double d = 3.14;
char c = 'A';
bool flag = true;

double j = double(i);        // or (double) i
int e = int(d);              // or (int) i
int d = int(c);              // or (int) c
int flag2 = bool(flag);      // or (int) flag

// ...
\end{console}
All the above can be executed using \textbf{\texttt{static\_cast}}:

\begin{console}
int i = 42;
double d = 3.14;
char c = 'A';
bool flag = true;

double i1 = static_cast< double >(i);
int d1 = static_cast< int >(d);
int c1 = static_cast< int >(c);
bool flag1 = static_cast< bool> (flag);

// ...
\end{console}

In general static cast is allowed when automatic typecasting is allowed.
For instance

\begin{console}
int i = 42;
double i1 = i;                        // automatic
double i2 = double(i);                // OK
double i3 = static_cast< double >(i); // OK
\end{console}

However this is not allowed:

\begin{console}
int i = 42;
double * p = &i;
\end{console}

Therefore this will not compile:

\begin{console}
int i = 42;
double * p = static_cast< double *>(&i); // BAD!!!
\end{console}

Of course C-style casting allows you to do this:

\begin{console}
int i = 42;
double * p = (double *)(&i);
\end{console}

And this is \EMPHASIZE{the whole point of C++ typecasting}. It's to subdivide C style castings into various kinds of C++ type casting so as to provide finer control and to provide more warnings.

That being said, note that in many areas of C++ programming, as long as you know what you are doing, i.e., you are only typecasting when automatic typecasting is allowed (and not crazy typecasting between pointers of different types), then you'll see that C-style typecasting is still used when writing C++ code.

For instance when you look OpenCV programs (CV = computer vision, OpenCV is a very famous open source CV library written in C/C++), it's very common to see C-style typecasting rather than C++ style typecasting. Also, OpenGL C++ program (for computer graphics) tend to use C-style typecasting.

